\chapter{דו־כיווניות בפרסום אקדמי}

\section{מהי דו־כיווניות ושפות \textenglish{RTL}}

דו־כיווניות, או \textenglish{BiDi}, היא יכולת של מערכות צפייה וניתוח לטפל בשפות
הנכתבות משמאל לימין (כמו עברית וערבית) וגם בשפות הנכתבות מימין לשמאל (כמו אנגלית וצרפתית)
באותו מסמך. זו בעיה מורכבת מפני שסדר הקריאה של הטקסט שונה משדר הטקסט הלוגי.

\section{אתגרי דו־כיווניות בטיפוגרפיה}

כאשר משלבים עברית עם אנגלית בנוסח אקדמי, יש צורך להתמודד עם כמה אתגרים:

\begin{equation}
\text{BiDi Level} = f(\text{language}, \text{script}, \text{punctuation})
\end{equation}

דוגמה פשוטה: כאשר כותבים את שם הדגם \textenglish{BERT} באמצע משפט עברית,
פרקים כמו סוגריים ומקפים עלולים להתהפך במידה לא נכונה.

\section{פתרון באמצעות \textenglish{polyglossia} ו־\textenglish{bidi}}

בעת שימוש ב־\textenglish{LuaLaTeX} עם \textenglish{polyglossia}, 
המערכת אוטומטית טוענת את חבילת ה־\textenglish{bidi}.
אנו מגדירים את העברית כשפה ראשית ואת האנגלית כשפה משנית בפרמבולה:

\begin{LTR}
\begin{verbatim}
\setmainlanguage{hebrew}
\setotherlanguage{english}
\end{verbatim}
\end{LTR}

זה אומר ללוח הרינדור להתקן את הטקסט מימין לשמאל כברירת מחדל,
וכן לעדכן את כיוון הקריאה במקומות מסוימים בהם מופיע טקסט אנגלית.

\section{השימוש ב־\textenglish{\textbackslash textenglish\{\}}}

עבור מילים או ביטויים אנגליים בודדים בתוך טקסט עברי, אנו משתמשים בפקודה \textenglish{\textbackslash textenglish\{\}}.
זה מדגיש שהטקסט בתוך הסוגריים הוא אנגלית ולא עברית, ולפיכך לא צריך להיות בכיוון \textenglish{RTL}:

\begin{equation}
\text{עברית} \rightarrow \textenglish{English} \rightarrow \text{עברית}
\end{equation}

לדוגמה, כאשר אנו מתייחסים למודלים כמו \textenglish{Transformer}, \textenglish{GPT}, ו־\textenglish{BERT},
כל שם מודל חייב להיות מעטופף ב־\textenglish{\textbackslash textenglish\{\}}.

\section{בלוקים של קוד וסימון מתמטי}

עבור בלוקים מרובי־שורות של קוד או נוסחאות מתמטיות שחייבות להיות משמאל לימין,
אנו משתמשים בסביבת \textenglish{\textbackslash begin\{LTR\}...\textbackslash end\{LTR\}}:

\begin{LTR}
\begin{verbatim}
def attention(Q, K, V):
    scores = (Q @ K.T) / sqrt(d_k)
    weights = softmax(scores)
    return weights @ V
\end{verbatim}
\end{LTR}

הסביבה \textenglish{LTR} מוודא שסימן פיתוח כמו \textenglish{(} או \textenglish{[}
לא מתהפך לסימן סגור כמו \textenglish{)} או \textenglish{]}.

\section{בדיקת המערכת}

כדי לבדוק שכל זה עובד כראוי, יש להריץ:

\begin{LTR}
\begin{verbatim}
lualatex -interaction=nonstopmode main.tex
\end{verbatim}
\end{LTR}

אם אין שגיאות ותהליך הקומפילציה מסתיים בהצלחה,
אז כל בדיקות דו־כיווניות עברו בהצלחה.
